\documentclass[12pt, a4paper]{article}

% Γλώσσα και Κωδικοποίηση
\usepackage[utf8]{inputenc}
\usepackage[greek,english]{babel}
\selectlanguage{greek}
\usepackage{alphabeta}
\usepackage{textgreek}  % Για ελληνικά γράμματα σε μαθηματική λειτουργία
\newcommand{\lt}[1]{\textlatin{#1}}
\usepackage{enumitem}

% Μαθηματικά
\usepackage{amsmath}
\usepackage{amssymb}

% Μορφοποίηση Σελίδας και Διάστιχο
\usepackage{geometry}
\geometry{left=3cm, right=2cm, top=4cm, bottom=3cm}
\usepackage{setspace}
\setstretch{1.5}

% Χρώματα και Ορισμοί
\usepackage{xcolor}
\definecolor{maroon}{HTML}{AF3235}
\newcommand{\HRule}{\rule{\linewidth}{0.5mm}}

% Πακέτο για monospace κώδικα
\usepackage{verbatim}

% Εικόνες
\usepackage{graphicx}
\graphicspath{{images/}}

% Εντολή για άδεια σελίδα
\usepackage{afterpage}
\newcommand\blankpage{%
	\null
	\thispagestyle{empty}%
	\addtocounter{page}{-1}%
	\newpage}

\begin{document}
	
\begin{titlepage}
	\begin{center}
		\HRule \\[0.6cm] 
		{\huge \bfseries \strut Numerical Analysis - Lab Project 1} \\[0.4cm] 
		\HRule \\[1.5cm]
		{\large \textbf{Topic:}} \\
		{\large Decimal to Base-b Positional System Conversion Utility} \\[2cm]
		
		\vfill 
		\begin{doublespacing}
			{\huge Antonios Zacharias\\}
			{\large \textbf{BSc} in Mathematics, AUTh\\}
			{\large \textbf{MSc} in Communication Networks \& System Security \\
			Computer Science Department, AUTh}
		\end{doublespacing}
	
		\vspace*{1cm} 
		\end{center}
\end{titlepage}	
	
\section*{1. Integer Part Conversion (intdec2baseb)}
\textbf{Implementation Logic:} The function \texttt{intdec2baseb} converts the integer part of a decimal number to any target base $b \in [2, 9]$.

\begin{itemize}[label=--]
	\item \textbf{Algorithm:} It implements successive Euclidean divisions by the target base $b$.
	\item \textbf{Mechanism:} A \texttt{while} loop runs as long as the quotient is greater than zero.
	\item \textbf{Calculation:} The \texttt{mod(i, b)} function isolates the remainder, while \texttt{floor(i/b)} updates the integer for the next iteration.
	\item \textbf{Output:} Each resulting digit is converted to a string and concatenated using \texttt{strcat}.
\end{itemize}

\textbf{Example (Base-5):}
\begin{verbatim}
	>> intdec2baseb(126, 5)
	ans = '1001'
\end{verbatim}

\section*{2. Fractional Part Conversion (fractdec2baseb)}
\textbf{Implementation Logic:} The function \texttt{fractdec2baseb} handles the conversion of decimal fractions ($< 1$) to the target base $b$.

\begin{itemize}[label=--]
	\item \textbf{Precision Control:} A counter and the input parameter \texttt{ndigits} ensure the algorithm terminates after a user-defined number of digits.
	\item \textbf{Algorithm:} The decimal fraction is repeatedly multiplied by the target base $b$.
	\item \textbf{Digit Isolation:} The \texttt{floor} function extracts the integer part of each product, which represents the next fractional digit.
	\item \textbf{Formatting:} The resulting digits are joined into a string, prefixed with a decimal point.
\end{itemize}

\textbf{Example (Base-5):}
\begin{verbatim}
	>> fractdec2baseb(0.1234, 5, 10)
	ans = '.0302030303'
\end{verbatim}

\section*{3. Unified Base Conversion (mydec2baseb)}
\textbf{Implementation Logic:} The primary function \texttt{mydec2baseb} serves as a generalized wrapper for full decimal-to-base-$b$ conversion.

\begin{itemize}[label=--]
	\item \textbf{Integration:} It splits the input decimal number into its integer and fractional components using \texttt{floor}.
	\item \textbf{Modular Execution:} It calls \texttt{intdec2baseb} for the integer part and \texttt{fractdec2baseb} for the fractional part.
	\item \textbf{Final Assembly:} The two outputs are merged into a single string using \texttt{strcat}.
\end{itemize}

\textbf{Execution Examples:}
\begin{verbatim}
	>> mydec2baseb(98.5, 9, 7)
	ans = '118.4444444'
	
	>> mydec2baseb(139.25, 4, 8)
	ans = '2023.1'
	
	>> mydec2baseb(126.1234, 5, 10)
	ans = '1001.0302030303'
	
	(Examples Verified via MATLAB Command Window)
\end{verbatim}

\end{document}



